% ---------------------------------------------------------------------
%  CHƯƠNG 7 — KẾT LUẬN VÀ HƯỚNG PHÁT TRIỂN
% ---------------------------------------------------------------------
\chapter{Kết luận và hướng phát triển}
\label{ch:ketluan}

\begin{muctieu}
Chương này tổng kết các đóng góp chính của đồ án, rút ra các bài học kinh nghiệm trong quá trình xây dựng hệ thống AI gợi ý món ăn, và đề xuất các định hướng nghiên cứu, phát triển trong tương lai.
\end{muctieu}

\section{Kết luận}

Đồ án \en{``Hệ thống AI gợi ý món ăn theo nguyên liệu''} đã hoàn thành xuất sắc các mục tiêu đề ra ban đầu cả về mặt lý thuyết giải thuật lẫn hiện thực hoá ứng dụng thực tế:
\begin{enumerate}[leftmargin=1.6em, itemsep=2pt]
  \item \textbf{Đóng góp học thuật và giải thuật:} Đề xuất và cài đặt thành công mô hình gợi ý ẩm thực kết hợp Grounded Database-First RAG giữa mô hình ngôn ngữ lớn Google Gemini AI và kho tri thức ẩm thực quan hệ MySQL. Giải thuật triệt tiêu hoàn toàn vấn đề ảo giác nguyên liệu thịt cá lạ (đạt 100\% độ chuẩn xác loại trừ), đồng thời cung cấp bảng phân tích định lượng 8 chỉ số dinh dưỡng khoa học.
  \item \textbf{Đóng góp thực tiễn và sản phẩm:} Xây dựng hoàn chỉnh ứng dụng Web toàn diện (React 18 + Node.js Express + MySQL) với đầy đủ 2 phân hệ Khách hàng và Quản trị viên. Kho tri thức ẩm thực sở hữu 552 món ăn Việt Nam truyền thống và hơn 2.000 hình ảnh Cookpad CDN sắc nét.
  \item \textbf{Độ bền vững cao:} Kiến trúc chuyển mạch tự động Fallback bảo vệ hệ thống luôn hoạt động ổn định với thời gian phản hồi chỉ 0,21 giây (< 300ms) ngay cả khi mất kết nối mạng bên ngoài hoặc hết quota API.
  \item \textbf{Triển khai thực tế:} Hệ thống đã được đưa vào môi trường sản xuất thực tế trên Vercel Edge Network và Render Cloud, hoạt động ổn định với người dùng thật.
\end{enumerate}

\section{Bài học rút ra}

Qua quá trình thực hiện đồ án, nhóm tác giả đã tích lũy được nhiều bài học kinh nghiệm quý báu:
\begin{itemize}[leftmargin=1.6em, itemsep=2pt]
  \item \textbf{Kiểm soát mô hình AI tạo sinh:} Mô hình ngôn ngữ lớn (LLM) có năng lực sáng tạo rất cao nhưng không thể phó mặc hoàn toàn cho AI trong các bài toán đòi hỏi tính chính xác thực tế. Cần phải áp dụng các kỹ thuật định hướng ngữ cảnh như Grounded RAG, hạ thấp nhiệt độ sáng tạo (temperature) và thiết lập các lớp phòng thủ kiểm tra dữ liệu ở tầng backend.
  \item \textbf{Đặc thù tiếng Việt trong AI:} Xử lý ngôn ngữ tự nhiên tiếng Việt đòi hỏi phải chú trọng đến dấu thanh, hiện tượng từ ghép và ranh giới từ nguyên vẹn để tránh các lỗi so khớp nguy hiểm trong đời sống ẩm thực.
  \item \textbf{Tư duy thiết kế chịu lỗi:} Luôn phải xây dựng phương án dự phòng (Fallback) khi tích hợp các dịch vụ đám mây của bên thứ ba để đảm bảo hệ thống duy trì tính sẵn sàng cao (High Availability).
\end{itemize}

\section{Hướng phát triển}

Để hoàn thiện và nâng cao hơn nữa giá trị thực tiễn của sản phẩm, nhóm đề xuất một số định hướng phát triển trong tương lai:
\begin{itemize}[leftmargin=1.6em, itemsep=2pt]
  \item \textbf{Tích hợp Thị giác máy tính (Computer Vision):} Huấn luyện mô hình YOLOv8 hoặc YOLOv11 để người dùng chỉ cần chụp ảnh các ngăn thực phẩm trong tủ lạnh, hệ thống sẽ tự động phân loại và nhận diện danh sách nguyên liệu đưa vào động cơ gợi ý.
  \item \textbf{Lập kế hoạch thực đơn theo tuần (Weekly Meal Planner):} Mở rộng thuật toán cho phép người dùng nhập mục tiêu thể trạng (giảm cân, tăng cơ, ăn kiêng Keto, bệnh nhân tiểu đường) để AI tự động lên thực đơn 7 ngày cân đối calo và xuất danh sách đi chợ tiện lợi.
  \item \textbf{Phát triển ứng dụng di động đa nền tảng:} Xây dựng ứng dụng di động native bằng React Native hoặc Flutter để mang lại trải nghiệm tiện dụng tối đa cho người dùng ngay tại gian bếp.
  \item \textbf{Mở rộng kho tri thức ẩm thực đa vùng miền:} Thu thập thêm các món ăn đặc sản vùng miền của đồng bào các dân tộc thiểu số và ẩm thực quốc tế.
\end{itemize}
